\documentclass[../../../include/open-logic-section]{subfiles}

\begin{document}

\olfileid{sth}{ordinals}{replacement}
\olsection{替换公理}

在 \olref[sth][ordinals][vn]{sec} 中，我们曾建议将序数视为序型，即良序的典范代表，以此作为引入序数的动机。为此，我们需要证明\emph{每个良序都同构于某个序数}。这就允许我们将 $\ordtype{A, <}$ 定义为满足 $\tuple{A, <} \isomorphic \alpha$ 的序数 $\alpha$。

不幸的是，仅凭目前引入的公理，我们\emph{无法}证明这一期望的结果。（我们将在 \olref[replacement][strength]{sec} 中看到原因，但目前的关键是：我们做不到。）我们需要一个新的想法，如下：

\begin{axiom}[替换公理模式]
对于任意公式 $\phi(x, y)$，以下是一条公理：
\begin{quote}
	对于任意 $A$，如果 $(\forall x \in A)\lexists![y][\phi(x,y)]$，那么
	$\Setabs{y}{(\exists x \in A)\phi(x,y)}$ 存在。
\end{quote}
\end{axiom}

\noindent
与分离公理一样，这是一个模式：它产生无穷多条公理，对应于无穷多个不同的 $\phi$。它也可以（并且通常是）这样表述：

\begin{defish}
对于任意不包含“$B$”的公式 $\phi(x,y)$，以下是一条公理：
\[
\forall A[(\forall x \in A)\lexists![y][\phi(x,y)] \lif \exists B\forall y (y \in B \liff (\exists x \in A)\phi(x,y))]
\]
\end{defish}

然而，初看之下，这是一个相当复杂的公式。替换公理的以下直接推论，或许能更\emph{清晰}地表达我们所依据的直观想法：\footnote{项是在对其自由变元的任何代入下恰好指称一个对象的表达式。例如，“$\emptyset$”是指称空集的项；“$\{x\}$”是指称 $x$ 的单元素集的项（无论 $x$ 是什么）；“$x \cup y$”是指称 $x$ 和 $y$ 的并集的项（无论它们是什么）。}

\begin{cor}
对于任意项 $\tau(x)$ 和任意集合 $A$，以下集合存在：
\[
	\Setabs{\tau(x)}{x \in A} = \Setabs{y}{(\exists x \in A)y = \tau(x)}.
\]
\end{cor}

\begin{proof}
由于 $\tau$ 是一个\emph{项}，$\forall x \lexists![y][\tau(x) = y]$。
更显然地，$(\forall x \in A)\lexists![y][\tau(x) = y]$。因此，
由替换公理可知 $\Setabs{y}{(\exists x \in A)\tau(x) = y}$ 存在。
\end{proof}

\noindent
这表明“替换”是这个公理的一个好名字：给定集合 $A$，你可以用 $A$ 的每个成员在~$\tau$ 下的像替换该成员，从而形成新的集合 $\Setabs{\tau(x)}{x \in A}$。实际上，沿用函数下集合的像的记法，我们可以将 $\Setabs{\tau(x)}{x \in A}$ 写作 $\funimage{\tau}{A}$。

然而，关键在于 $\tau$ 是一个\emph{项}。它不必是 \olref[sfr][fun][rel]{sec} 意义下的\emph{函数}（即某个有序对的集合）的名字。毕竟，如果 $f$ 是该意义下的函数，那么集合 $\funimage{f}{A} = \Setabs{f(x)}{x \in
A}$ 仅仅是 $\ran{f}$ 的一个特定子集，而其存在性已由~$\Zminus$ 的公理所保证。\footnote{只需考虑 $\Setabs{y \in \bigcup \bigcup f}{(\exists x \in A)y =
f(x)}$。} 相比之下，替换公理是对我们公理体系的一个\emph{强大}补充，正如我们将在 \olref[sth][replacement][]{chap} 中看到的那样。

\end{document}